\begin{table*}[ht]
\centering
\caption{Requirements for the framework}
\label{tab:reliability-architecture-requirements}
\footnotesize
\begin{tabular}{@{}p{0.8cm}p{10cm}p{3.5cm}@{}}
\toprule
\textbf{ID} & \textbf{Requirement} & \textbf{Synthesized From} \\ 
\midrule

\textbf{R1} & Provide a \textbf{compensation mechanism} to handle missing or inconsistent event inputs through reconstruction or substitution. & \citet{avizienis2004basic}, \citet{cugola2015complex}, \chref{ch:literature-review} \\

\textbf{R2} & Enable \textbf{confidence modeling and propagation} so that events include quantitative confidence values maintained through the CEP pipeline. & \citet{cugola2015complex}, \citet{wasserkrug2006taxonomy}, \citet{hanmer2013patterns}, \secref{subsec:background-cep-uncertainty} \\

\textbf{R3} & Support \textbf{graceful degradation} by continuing partial event processing under uncertainty or data loss. & \citet{hanmer2013patterns}, \secref{subsec:literature-review-results-rq5} \\

\textbf{R4} & Maintain a \textbf{modular and extensible architecture} composed of loosely coupled, replaceable components. & \secref{para:lt-rq2-sots-type}, \figref{fig:sos-dim} \\

\textbf{R5} & Provide \textbf{runtime observability} through measurable indicators that reflect reliability and compensation performance. & \secref{subsec:literature-review-results-rq5} \\

\textbf{R6} & Ensure \textbf{interoperability} through standardized data formats and lightweight communication protocols. & \secref{para:lt-rq1-challenges} \\

\bottomrule
\end{tabular}
\end{table*}
